In $\mathbb{Q}$ sono definite le operazioni di somma e prodotto nel seguente modo.

\begin{definition}
    \textbf{Somma in $\mathbb{Q}$}
    \[
        \forall r_1 = \left[\frac{a}{b}\right]_\sim, r_2 = \left[\frac{c}{d}\right]_\sim \in \mathbb{Q}, \;
        r_1 + r_2 := \left[\frac{ad+bc}{bd}\right]
    \]
\end{definition}
\begin{definition}
    \textbf{Prodotto in $\mathbb{Q}$}
    \[
        \forall r_1 = \left[\frac{a}{b}\right]_\sim, r_2 = \left[\frac{c}{d}\right]_\sim \in \mathbb{Q}, \;
        r_1 \cdot r_2 := \left[\frac{ac}{bd}\right]
    \]
\end{definition}

L'operazione di somma in $\mathbb{Q}$ gode delle seguenti proprietà.

\begin{property}
    \textbf{Commutativa}
    \[\forall a, b \in \mathbb{Q}, \, a + b = b + a.\]
\end{property}
\begin{property}
    \textbf{Associativa}
    \[\forall a, b, c \in \mathbb{Q}, \, (a + b) + c = a + (b + c).\]
\end{property}
\begin{property}
    \textbf{Esistenza dell'elemento neutro}
    \[\forall a \in \mathbb{Q}, \, a + 0 = 0 + a = a.\]
\end{property}
\begin{property}
    \textbf{Esistenza dell'elemento inverso}
    \[\forall a \in \mathbb{Q}, \, \exists b = -a \in \mathbb{Q} \mid a + b = 0.\]
\end{property}
\begin{property}
    \label{p:raz-chi-som}
    \textbf{Chiusura di $\mathbb{Q}$ rispetto alla somma}
    \[\forall a, b \in \mathbb{Q}, \; (a + b) \in \mathbb{Q}.\]
\end{property}
La struttura algebrica $(\mathbb{Q}, +)$ è dunque un gruppo abeliano.
\vspace{0.4cm}

L'operazione di prodotto in $\mathbb{Q}' := \mathbb{Q} \setminus \left\{0\right\}$ gode delle seguenti proprietà.

\begin{property}
    \textbf{Commutativa}
    \[\forall a, b \in \mathbb{Q}', \, a \cdot b = b \cdot a.\]
\end{property}
\begin{property}
    \textbf{Associativa}
    \[\forall a, b, c \in \mathbb{Q}', \, (a \cdot b) \cdot c = a \cdot (b \cdot c).\]
\end{property}
\begin{property}
    \textbf{Esistenza dell'elemento neutro}
    \[\forall a \in \mathbb{Q}', \, a \cdot 1 = 1 \cdot a = a.\]
\end{property}
\begin{property}
    \label{p:inv-q}
    \textbf{Esistenza dell'elemento inverso}
    \[\forall a \in \mathbb{Q}', \, \exists b = a^{-1} = \frac{1}{a} \in \mathbb{Q}' \mid a \cdot a^{-1} = a \cdot \frac{1}{a} = 1.\]
\end{property}
\begin{property}
    \textbf{Chiusura di $\mathbb{Q}$ e $\mathbb{Q}'$ rispetto al prodotto}
    \label{p:raz-chi-pro}
    \[\forall a, b \in \mathbb{Q}, \; (a \cdot b) \in \mathbb{Q},\]
    \[\forall a, b \in \mathbb{Q}', \; (a \cdot b) \in \mathbb{Q}'.\]
\end{property}
La struttura algebrica $(\mathbb{Q} \setminus \left\{0\right\}, \cdot)$ è dunque un gruppo abeliano.

\begin{observation}
    Viene escluso l'elemento 0 poiché altrimenti non sarebbe garantita \hyperref[p:inv-q]{\textbf{p~\ref*{p:inv-q}}}
\end{observation}


\vspace{0.4cm}

In $\mathbb{Q}$ vale inoltre la proprietà distributiva del prodotto sulla somma
\begin{property}
    \textbf{Distribtiva di $\cdot$ su +}
    \[\forall a, b, c \in \mathbb{Q}, \; a \cdot (b + c) = a \cdot b + a \cdot c.\]
\end{property}